La diversidad de equipos y configuraciones puede desplazar tiempo de formación científica hacia la resolución de problemas de instalación. Se describe una arquitectura de entornos de trabajo reproducibles basada en repositorios Git, contenedores de desarrollo, cuadernos y automatización, contextualizada en SILAFAE XV y la CERN School of Computing de América Latina de 2026. El enfoque conecta preparación del entorno, ejecución del análisis y conservación de resultados. Se presenta la observación comunicada por el autor de una preparación para 80 estudiantes en menos de diez minutos, indicando que su interpretación requiere precisar el intervalo y recuperar registros operativos. Se distingue esa observación de una evaluación controlada del aprendizaje. Como demostrador complementario, los propios manuscritos se gestionan en LaTeX y generan PDF mediante integración continua. Se discuten conectividad, cuotas, mantenimiento, portabilidad y una propuesta de evaluación reproducible para futuras escuelas.
